\section{Background}
An IOMMU associates a device request with a device context, translates its I/O virtual address, and checks permissions before forwarding the request. The IOTLB caches translation information. On a miss, a page-table walker (PTW) traverses the applicable page tables; device-directory or other data-structure accesses can also be necessary. The RISC-V specification defines the architectural behavior and management interfaces, while the organization and parallelism of the implementation determine how requests share service~\cite{riscvIommu}.\claim{B-SPEC}

Three forms of delay matter to this characterization. \emph{Shared data-path competition} exists when concurrent DMA reads, writes, and responses use common downstream resources. \emph{Service serialization} makes one translation wait for another request even when its eventual lookup hits. \emph{Capacity replacement} removes a reusable translation, causing a later request to miss.\claim{B-MECHANISMS} These mechanisms can coexist, so a miss counter is neither a decomposition of total delay nor evidence that every recorded hit received immediate service.

The distinction has substantial precedent. HERO exposes translation events and demonstrates hit-under-miss behavior. Kurth et al. explicitly describe how full miss-related buffers can stall later accesses, including hits, and examine scalable handling and DMA support~\cite{hero2017,kurth2018}.\claim{L-PREHISTORY} Our local intervention tests the corresponding dependency in a different fixed controller, rather than claiming a new general principle.

We call the fixed foreground stream the \emph{victim} (V) and its co-running stream the \emph{aggressor} (A). These names identify experimental roles, not malicious intent. Solo runs V alone; A1 provides a co-running small-working-set control, while A8, A16, and A64 increase the aggressor's configured working set.\claim{M-CONTROLS} Equal nominal transfer sizes do not imply equal translation activity or equal service schedules.
